\chapter{Configure}\label{ch:configure}

\emph{You want Agern reachable under its proper address, ideally over HTTPS, protected by a
firewall, and open to the other people on your team.} This chapter covers the few settings a new
installation really touches. Every setting is described in
\repodoc{ARCHITECTURE.md\#3-configuration-environment-variables}{\texttt{ARCHITECTURE.md}, section~3};
appendix~\ref{app:config} lists the ones you are likely to need.

After every change to \file{/etc/agern/agern.env}, restart Agern so that it reads the file again:
\begin{lstlisting}
sudo agern restart
\end{lstlisting}

\section{The public URL}\label{sec:public-url}

\envvar{FACTORY_PUBLIC_URL} is the one address under which people use Agern. The browser sends
it as the \emph{origin} of every request, and Agern refuses changes and live connections that come
from any other origin. This protects you from other web pages that try to act in your name.

\begin{enumerate}
  \item Decide on the address people will type, for example \url{http://agern.example.com:8080}
    on your LAN, or \url{https://agern.example.com} behind a reverse proxy (next section).
  \item Set it, without a trailing slash, and restart:
\begin{lstlisting}[style=conf]
FACTORY_PUBLIC_URL=https://agern.example.com
\end{lstlisting}
  \item Open exactly that address. If you use another one (the IP address, \texttt{localhost}, a
    different port), pages may load, but signing in or anything else that changes data fails with
    \msg{Origin not allowed}.
\end{enumerate}

\section{HTTPS behind a reverse proxy}\label{sec:https}

Plain HTTP works on a trusted LAN: the second factor is a code from an app, which does not need
HTTPS. For HTTPS, put a reverse proxy such as Caddy or nginx on the same machine in front of
Agern.

\begin{enumerate}
  \item Make Agern listen only on the loopback address, so that it is reachable only through the
    proxy, and set the HTTPS address as public URL. An \texttt{https://} URL also makes the session
    cookie \texttt{Secure}:
\begin{lstlisting}[style=conf]
FACTORY_HOST=127.0.0.1
FACTORY_PUBLIC_URL=https://agern.example.com
\end{lstlisting}
  \item Configure the proxy to forward everything to \texttt{127.0.0.1:8080}. It must pass
    WebSockets (Agern streams agents live over \texttt{/api/ws/\ldots}) and should send the header
    \texttt{X-Forwarded-For}. Caddy does both by default; for a LAN name, \texttt{tls internal}
    makes Caddy use its own certificate authority (\file{/etc/caddy/Caddyfile}):
\begin{lstlisting}[style=conf]
agern.example.com {
    tls internal
    reverse_proxy 127.0.0.1:8080
}
\end{lstlisting}
    With nginx, add WebSocket upgrades and
    \texttt{proxy\_set\_header X-Forwarded-For \$proxy\_add\_x\_forwarded\_for;}.
  \item Restart Agern and the proxy, then open \url{https://agern.example.com}.
\end{enumerate}

The forwarded header is more than logging: it is how Agern tells requests through the proxy from
requests made on the machine itself, which matters for automatic upgrades
(chapter~\ref{ch:upgrade}). Details:
\repodoc{INSTALL-LINUX.md\#7-network-firewall-and-tls}{\texttt{INSTALL-LINUX.md}, section~7}.

\section{The firewall}\label{sec:firewall}

Agern is made for a trusted network. Do not expose it to the internet.

\begin{enumerate}
  \item Allow SSH and Agern's port only from your LAN (here \texttt{192.0.2.0/24}; use port 443
    instead of 8080 behind an HTTPS proxy):
\begin{lstlisting}
sudo apt install ufw
sudo ufw default deny incoming
sudo ufw allow from 192.0.2.0/24 to any port 22 proto tcp
sudo ufw allow from 192.0.2.0/24 to any port 8080 proto tcp
sudo ufw enable
\end{lstlisting}
\end{enumerate}

\section{The essential settings}\label{sec:essential-settings}

Most installations change nothing beyond chapter~\ref{ch:install} and the address above. These
are the settings worth knowing about; their defaults suit a small team:

\begin{center}
\begin{tabularx}{\linewidth}{@{}>{\raggedright\arraybackslash}p{0.42\linewidth}>{\raggedright\arraybackslash}X@{}}
  \toprule
  Variable (default) & What it decides\\
  \midrule
  \envvar{FACTORY_HOST} (\texttt{0.0.0.0}), \envvar{FACTORY_PORT} (\texttt{8080}) & Where Agern listens.\\
  \envvar{FACTORY_MAX_AGENTS_PER_USER} (\texttt{8}) & How many agents one person may run at once.\\
  \envvar{FACTORY_STAGE_TIMEOUT_MINUTES} (\texttt{60}) & The longest a single pipeline stage may take before the run fails.\\
  \envvar{FACTORY_CHAT_IDLE_MINUTES} (\texttt{30}) & After how long an idle chat's process is stopped; your next message resumes it.\\
  \bottomrule
\end{tabularx}
\end{center}

\section{Adding users}\label{sec:users}

Agern has no user database of its own to manage: whoever can sign in with an account of the
machine (\texttt{os}) or of your directory (\texttt{ldap}) can use it. What a person then sees is
decided by their own GitHub token.

\begin{enumerate}
  \item \textbf{Accounts on this machine} (\texttt{FACTORY\_AUTH=os}): create one per person with
    \cmd{sudo adduser <name>}, with a password they use nowhere else. The \texttt{shadow} group
    decision of section~\ref{sec:first-admin} applies.
  \item \textbf{Or a directory} (\texttt{FACTORY\_AUTH=ldap}): set \envvar{FACTORY_LDAP_URL} and
    either a DN template or search-then-bind, for example:
\begin{lstlisting}[style=conf]
FACTORY_AUTH=ldap
FACTORY_LDAP_URL=ldaps://ldap.example.com:636
FACTORY_LDAP_USER_DN_TEMPLATE=uid={0},ou=people,dc=example,dc=com
\end{lstlisting}
    For search-then-bind, see the \texttt{FACTORY\_LDAP\_\ldots} lines in
    \file{/etc/agern/agern.env}. \texttt{FACTORY\_AUTH=os,ldap} tries both, in that order.
  \item \textbf{Admins}: add the user name to \envvar{FACTORY_ADMINS} and restart. Admins see the
    \ui{Admin} page, approve sandboxes and can stop or cancel anybody's agents, but never drive
    them.
  \item Each new person signs in, enrols an authenticator and saves their recovery codes, exactly
    as in section~\ref{sec:first-login}, then connects their own accounts
    (chapter~\ref{ch:accounts}).
\end{enumerate}

The \ui{Admin} page lists everybody who has signed in, whether they have a GitHub token and a
Claude credential, and their two-factor state (figure~\ref{fig:admin-users}). Someone who lost
their phone and their recovery codes, or locked themselves out with five wrong codes, needs an
admin:

\begin{enumerate}
  \item Confirm the person's identity outside Agern.
  \item On the \ui{Admin} page, choose \ui{Reset 2FA} in their row and confirm
    (figure~\ref{fig:admin-reset}). They enrol a new authenticator at their next sign-in.
\end{enumerate}

\screenshot[0.85\linewidth]{45-admin-users-desktop}{The users on the Admin page, with their credentials and two-factor state.}{admin-users}

\screenshot{22-admin-reset-2fa-desktop}{Resetting a user's two-factor setup.}{admin-reset}

\whatyouhave
Agern answers under its public address, over HTTPS if you chose to, only your LAN can reach it,
and your colleagues can sign in. Next, connect your GitHub and Claude accounts so that you see
projects and can start agents: chapter~\ref{ch:accounts}.
